\documentclass[11pt,a4paper]{article}

\usepackage[T1]{fontenc}
\usepackage[utf8]{inputenc}
\usepackage{lmodern}
\usepackage[a4paper,margin=26mm]{geometry}
\usepackage{amsmath,amssymb,amsthm,mathtools}
\usepackage{booktabs,array}
\usepackage{microtype}
\usepackage{xcolor}
\usepackage{enumitem}
\usepackage{url}
\usepackage[colorlinks=true,linkcolor=blue!55!black,citecolor=blue!55!black,
            urlcolor=blue!55!black]{hyperref}

\hypersetup{
  pdftitle={Order-Dependent Rejection Sampling in Rhyme-128},
  pdfauthor={Martin Feussner}
}

\newtheorem{theorem}{Theorem}
\newtheorem{lemma}[theorem]{Lemma}
\newtheorem{proposition}[theorem]{Proposition}
\newtheorem{remark}[theorem]{Remark}
\newcommand{\R}{\mathcal{R}}
\newcommand{\Z}{\mathbb{Z}}
\newcommand{\E}{\mathbb{E}}
\newcommand{\Sign}{\mathsf{Sign}}
\newcommand{\Verify}{\mathsf{Verify}}
\newcommand{\SampleChallenge}{\mathsf{SampleChallenge}}
\newcommand{\SampleInBall}{\mathsf{SampleInBall}}
\newcommand{\SHAKE}{\mathsf{SHAKE256}}
\newcommand{\pk}{\mathsf{pk}}
\newcommand{\sk}{\mathsf{sk}}
\newcommand{\norm}[1]{\left\lVert #1\right\rVert}
\newcommand{\infnorm}[1]{\left\lVert #1\right\rVert_{\infty}}
\newcommand{\trans}{\mathsf{T}}

\title{\bfseries Order-Dependent Rejection Sampling in Rhyme-128:\\[0.25em]
\large Exact Secret-Tail Recovery and a Fresh-Message Forgery\\
Using at Most 40,000 Signatures}
\author{Martin Feussner\\
\small Selmer Center, University of Bergen\\
\small \texttt{martin.feussner@uib.no}}
\date{1 October 2026}

\begin{document}
\maketitle

\begin{center}
\fbox{\parbox{0.92\linewidth}{\small\textbf{Provenance and review status.}
The attack was independently found by an AI system, \mbox{OpenAI Codex} using
Daybreak Blue at maximum reasoning effort, without a human first proposing
this specific attack.  The AI also produced the derivation, implementation,
experiments, adversarial review, and initial manuscript.  At the date above,
no human cryptanalyst has independently verified the argument,
implementation, measurements, or novelty assessment.  This preliminary
disclosure is provided for human review and reproduction.}}
\end{center}
\vspace{0.7em}

\begin{abstract}
Rhyme is a lattice signature submitted to the Next-generation Commercial
Cryptographic Algorithms Program.  Its coefficient rejection sampler scans a
parity-constrained set ``in a fixed order'', but the specification does not
define that order.  Appendix~A states the correct cumulative-envelope
requirement for its normalization constant, but then claims that the maximum
occurs at the boundary and prescribes the resulting single-term closed form.
For the operational Rhyme-128 CDT/source convention, exhaustive evaluation
gives prescribed values $9.3573397\cdot10^{-6}$ and
$1.4950374\cdot10^{-6}$ for the two
challenge parities, while the required maxima are respectively
$1.0019364869$ and $1.0019292771$.  Reading the PDF's displayed Gaussian
definition and parameter table literally gives different numbers, but the
same defect: boundary values $0.0075371423$ and $0.0035191981$ versus required
maxima $1.1818109123$ and $1.1821799192$.  In the operational completion
tested here, combining the prescribed boundary values with a deterministic
central-first order permitted by Algorithm~3 makes the sampled offset almost
equal to the negative challenge.  Each signature then gives a noisy linear
observation of the four-polynomial secret tail.

We recover the exact 1,024 secret-tail coefficients by negacyclic
least-squares regression and use the recovered public-key relation to forge at
zero commitment.  Two independent Rhyme-SHAKE-128 keys were tested on a
four-core Intel Core i5-6500 with 16~GiB of RAM.  The specification-tape run
recovered at 40,000 signatures; an independent-rejection-randomness run
recovered at 30,000.  Both used a 50,000-signature collection, with the final
10,000 held out for diagnostics.  An untouched submitted verifier accepted
and decoded all 50,000 attack-oracle signatures in each run.  Collection took
178.30 and 160.44 seconds, and the corresponding 40,000-signature batch
recoveries took 9.87 and 6.97 seconds on the shared machine.  In both runs,
recovery and forgery without access to the target secret produced
a 1,372-byte signature on a fresh message that the pristine submitted
verifier accepted.  Wrong-message, perturbed-tail, and wrong-key controls
rejected.

The result has a narrow but security-critical scope.  It is a practical
chosen-message attack on a deterministic completion that uses the fixed order
permitted by Algorithm~3 and the boundary value prescribed by Appendix~A.
The submitted source uses an ascending scan and sets
$M=1.002214572\ldots$, which satisfies the envelope for the operational CDT
distribution; that sampler is not broken by this attack.  A paired
50,000-signature control used the same key, messages, central order, and
signing logic as Run~2, changing only $M$ to that operationally safe value.
It failed the public-key check and recovered only 399 of 1,024 coefficients.
The evidence establishes a specification/design-level defect and an
internally inconsistent submitted specification.  It does not establish that
every implementation of the underlying abstract construction is broken.
\end{abstract}

\section{Introduction}

Fiat--Shamir-with-aborts signatures publish a response obtained by combining
a random mask, a challenge, and a long-term secret.  Rejection sampling is
used to make the accepted response distribution independent of that secret.
The normalization in such a sampler must dominate the complete acceptance
mass for every possible mask.  If it does not, a sequential implementation
can cross probability one before scanning every candidate, and an otherwise
irrelevant enumeration order becomes security-critical.

Rhyme uses a scalar rejection procedure for the first response component
\cite{rhyme-spec}.  Algorithm~3 assigns a probability increment to each
element of a finite parity set and scans the set ``in a fixed order.''  The
order is left unspecified.  Appendix~A first states the correct cumulative
envelope, but later asserts that a boundary point dominates its maximum and
prescribes the single-term expression obtained there.  Direct evaluation
shows that this expression is five to six orders of magnitude smaller than
the required cumulative envelope under the operational CDT convention and
more than two orders of magnitude smaller under a literal reading of the
PDF's Gaussian convention.

Consider the fixed central-minus order
\begin{align*}
 V_0&=(0,-2,2,-4,4,\ldots,-12,12),\\
 V_1&=(-1,1,-3,3,\ldots,-13,13).
\end{align*}
This is a deterministic permutation of exactly the specified sets, fixed
before the key, message, challenge, and signing randomness are known.  With
the prescribed boundary constant, an early central term normally exhausts
the uniform threshold.  The returned offset is therefore close to $v=-c$,
where $c$ is
the public binary challenge polynomial.  Rhyme's bottom response then has the
form
\[
 z_{\mathrm{bot}}\approx y_{\mathrm{bot}}-s_{\mathrm{tail}}c.
\]
The public challenges become regressors for the secret tail.

Our contributions are:
\begin{itemize}[leftmargin=1.5em]
  \item We show under both the literal PDF convention and the operational
        CDT/source convention that Appendix~A's boundary closed form violates
        the exact cumulative-envelope condition stated earlier in that
        appendix.
  \item We identify a second proof error concerning shifted discrete Gaussian
        sums and delimit its role in the failure.
  \item We instantiate a fixed central-first order allowed by the submitted
        specification and derive a public negacyclic least-squares recovery
        algorithm.
  \item We give an exact zero-commitment forgery from any recovered short tail
        satisfying the public-key relation.
  \item We demonstrate exact recovery and pristine-verifier-accepted
        fresh-message forgeries for two independent Rhyme-SHAKE-128 keys on
        the machine identified in Section~\ref{sec:experiment}.
  \item We separate the specification defect from the submitted source repair
        and confirm on the same key that the regression fails after the
        operationally valid envelope is restored.
\end{itemize}

The recovered object is the exact four-polynomial secret tail used in the
public relation, not the original 32-byte seed or the complete serialized
signing key.  The tail is nevertheless sufficient to produce valid signatures
on arbitrary fresh messages through the zero-commitment construction in
Section~\ref{sec:forgery}.  We therefore use the terms \emph{exact
secret-tail recovery} and \emph{fresh-message forgery} throughout.

\paragraph{Related public findings.}
As of 1 October 2026, the public Rhyme tracker lists three separate
findings: two high-level security ceilings in the submitted 384- and 512-bit
implementations and a missing doubled-width parity mask in the SM3 signers
\cite{ngcc-rhyme}.  The third finding gives noisy parity equations but does
not demonstrate key recovery or forgery.  We found no public report of the
invalid cumulative envelope, order-dependent sampler, or the attack described
here.  This is a best-effort novelty check, not a claim about unpublished or
unindexed work.

\section{Rhyme-SHAKE-128}
\label{sec:scheme}

Let
\[
  \R=\Z[x]/(x^{256}+1),\qquad q=3329.
\]
At Level~128, Rhyme uses $(k,\ell)=(2,3)$ and
$d=k+\ell-1=4$.  Write the short secret relation as
\[
                  s=(1,s_{\mathrm{tail}}),
       \qquad s_{\mathrm{tail}}=(s_1,s_2,s_3,s_4)\in\R^4.
\]
The public verification matrix is
\begin{equation}
 A=\bigl(-2b+qj\mid 2A_{\mathrm{gen}}\mid 2I_k\bigr)
       \in \R_{2q}^{k\times(k+\ell)},                 \label{eq:A}
\end{equation}
where $j$ selects the first public row.  Key generation ensures
\begin{equation}
                         As=qj\pmod {2q}.              \label{eq:keyrel}
\end{equation}
Each coefficient of the recovered tail lies in $[-2,2]$ in the submitted
parameter set.

For one signing attempt, Rhyme samples a top masking polynomial $y_1$ and
constructs
\begin{equation}
             y_{\mathrm{bot}}=2B_{\mathrm{sig}}X'+e_{\mathrm{bot}}.
                                                               \label{eq:ybot}
\end{equation}
Writing $\pk=(\mathrm{seed}_A,b_{\mathrm{gen}})$, Algorithm~6 computes
\begin{equation}
 \mu=H_{\mathrm{gen}}(\mathrm{seed}_A,b_{\mathrm{gen}},M),\qquad
 w=Ay,\qquad c=\SampleInBall(w,\mu),                 \label{eq:specchallenge}
\end{equation}
where $c$ is binary with Hamming weight $\tau=30$.  The submitted
Rhyme-SHAKE-128 source realizes Equation~\eqref{eq:specchallenge} by computing
\begin{align}
 \mu&=\SHAKE(\operatorname{enc}(\pk)\mathbin\Vert M)
                 [0\mathbin{:}\mathsf{CRHBYTES}],\nonumber\\
 \widetilde c&=\SHAKE(\operatorname{pack}_w(w)\mathbin\Vert\mu)
                 [0\mathbin{:}\mathsf{CTILDEBYTES}],\nonumber\\
 c&=\SampleChallenge(\widetilde c),                  \label{eq:sourcechallenge}
\end{align}
and transmits $\widetilde c$ together with the encoded response.  Here
$\operatorname{enc}(\pk)$ is the submitted byte encoding of
$(\mathrm{seed}_A,b_{\mathrm{gen}})$, and $[0\mathbin{:}L]$ denotes the first
$L$ output bytes.  Algorithm~3 acts on every coefficient pair
$(y_{1,i},c_i)$ and returns
\[
                           v_i=2x_i+c_i.
\]
With $v=(v_i)_i$, the transmitted response is
\begin{equation}
 z_1=y_1+v,\qquad
 z_{\mathrm{bot}}=y_{\mathrm{bot}}+s_{\mathrm{tail}}v.
                                                               \label{eq:response}
\end{equation}
All products are negacyclic in $\R$.  The verifier checks
\begin{equation}
                         w'=Az-qjc\pmod {2q}            \label{eq:verify}
\end{equation}
and Algorithm~7 recomputes $\mu$ and checks
$c=\SampleInBall(w',\mu)$.  The submitted source performs the corresponding
bytewise $\widetilde c$ comparison from Equation~\eqref{eq:sourcechallenge}.
At Level~128 the verifier also requires
\begin{equation}
 \infnorm{z_1}\le388,\qquad
 \infnorm{z_{\mathrm{bot}}}\le581,\qquad
 \norm{z}_2^2\le21{,}949{,}787.                       \label{eq:bounds}
\end{equation}

The serialized secret key contains more material than
$s_{\mathrm{tail}}$, including the signing basis and a seed $K$.  Our attack
does not recover those bytes and does not invoke the ordinary signing
algorithm after recovery.  It recovers the actual tail in
Equation~\eqref{eq:keyrel}, checks that relation publicly, and uses it to
construct signatures directly.

\section{The rejection-sampling defect}
\label{sec:defect}

For challenge bit $c\in\{0,1\}$, define
\begin{equation}
 V_c=\{v\in\Z:|v|\le R,\ v\equiv c\pmod 2\},\qquad
 p_c(v)=\frac{\rho_{\sigma_{\mathrm{tar}}}(v)}
               {\sum_{u\in V_c}\rho_{\sigma_{\mathrm{tar}}}(u)}.
                                                               \label{eq:pc}
\end{equation}
Rhyme-128 uses $R=13$ and a first-component bound $B_1=388$.  For a mask
$y\in[-B_1-R,B_1+R]$, Algorithm~3 scans $V_c$ in a fixed order.  If
$|y+v|\le B_1$, it adds
\begin{equation}
 a_y(v)=p_c(v)\frac{\rho_{\sigma_y}(y+v)}
                       {M_c\rho_{\sigma_y}(y)}         \label{eq:increment}
\end{equation}
to a cumulative value and returns the first $v$ for which that value exceeds
a uniform $r\in(0,1)$.

\begin{lemma}[Cumulative envelope]
\label{lem:envelope}
For a fixed challenge bit $c$, the sequential intervals in Algorithm~3 can
assign probability $a_y(v)$ to every eligible candidate only if
\begin{equation}
 M_c\ge
 \max_{|y|\le B_1+R}
 \sum_{\substack{v\in V_c\\|y+v|\le B_1}}
 p_c(v)\frac{\rho_{\sigma_y}(y+v)}{\rho_{\sigma_y}(y)}.
                                                               \label{eq:envelope}
\end{equation}
\end{lemma}

\begin{proof}
For fixed $(y,c)$, the intervals assigned to eligible candidates have total
length
\[
 \frac{1}{M_c}\sum_v p_c(v)
     \frac{\rho_{\sigma_y}(y+v)}{\rho_{\sigma_y}(y)}.
\]
If this exceeds one, later intervals lie partly or wholly outside the support
of $r$.  Their realized probabilities then depend on which candidates were
scanned first and cannot equal all of the prescribed increments.  Requiring
the total length to be at most one for every eligible $y$ and each parity
$c$ gives Equation~\eqref{eq:envelope}.
\end{proof}

Appendix~A of the submitted specification explicitly states
Equation~\eqref{eq:envelope}.  The defect is in its subsequent derivation.
After introducing an upper bound over all integers and completing the square,
the appendix treats
\begin{equation}
 \Theta_\alpha(\delta)=\sum_{v\in\Z}e^{-\alpha(v+\delta)^2}
                                                               \label{eq:theta}
\end{equation}
as translation invariant in $\delta$.  This is false for a noninteger shift.
Poisson summation gives
\[
 \Theta_\alpha(\delta)=\sqrt{\frac{\pi}{\alpha}}
 \sum_{k\in\Z}e^{-\pi^2k^2/\alpha}\cos(2\pi k\delta),
\]
which is periodic and is maximized at integer $\delta$, rather than constant.
Indeed,
\[
 \Theta_\alpha(0)-\Theta_\alpha(1/2)
 =4\sqrt{\frac{\pi}{\alpha}}
   \sum_{\substack{k\ge1\\k\ \mathrm{odd}}}e^{-\pi^2k^2/\alpha}>0.
\]
The difference is far below the displayed numerical precision for the
operational parameters.  Replacing the false equality by the valid inequality
$\Theta_\alpha(\delta)\le\Theta_\alpha(0)$ preserves the appendix's loose
universal upper-bound step, so this error alone does not produce the attack.

The security-critical step comes next.  Appendix~A evaluates the original
restricted sum where only one candidate is eligible at the extreme boundary,
asserts that this point dominates every interior $y$, and prescribes
\begin{equation}
 M_{\mathrm{bdry},c}=p_c(R_c)
  \frac{\rho_{\sigma_y}(B_1)}{\rho_{\sigma_y}(B_1+R_c)},
                                                               \label{eq:mbdry}
\end{equation}
where $R_c$ is the largest element of $V_c$.  This is one summand at one
boundary mask.  Exhaustive evaluation disproves the claimed dominance.

The displayed PDF formulas and the operational source instantiate different
Gaussian conventions.  Appendix~A literally defines
\begin{equation}
                  \rho_s(x)=e^{-\pi x^2/s^2}.           \label{eq:pdfgauss}
\end{equation}
For Rhyme-128, Table~3 lists $\sigma/(2\sqrt\pi)=1.295870$ and
$\sigma_y=122$, while the sampler text sets
$\sigma_{\mathrm{tar}}=2\sigma$.  Reading these statements literally gives
\[
 \sigma_{\mathrm{tar}}=4\sqrt\pi\,(1.295870)
       =9.187479087\ldots .
\]
The submitted source instead treats $122$ and $2.591740$ as ordinary standard
deviations.  Its operational weights are therefore
\begin{equation}
 \rho_y^{\mathrm{op}}(x)=e^{-x^2/(2\cdot122^2)},\qquad
 p_c^{\mathrm{op}}(v)\ \propto\ e^{-v^2/(2\cdot2.591740^2)},
                                                               \label{eq:opgauss}
\end{equation}
as confirmed independently by the CDT and Q63 tables.

Table~\ref{tab:envelope} evaluates Equation~\eqref{eq:envelope} under both
conventions, always retaining the submitted $R=13$ and $B_1=388$.  Every
integer $y\in[-401,401]$ and every eligible $v\in V_c$ was evaluated.  The
operational values were also recomputed directly from the source's fixed-point
Q63 probability table and Q62 exponential evaluator; they agree to every
digit displayed.

\begin{table}[t]
\centering
\small
\caption{Prescribed boundary constants and exhaustive cumulative-envelope
maxima under the two Gaussian conventions.}
\label{tab:envelope}
\begin{tabular}{@{}lcrrr@{}}
\toprule
Convention & $c$ & $M_{\mathrm{bdry},c}$ & $M_{\mathrm{req},c}$
    & $\operatorname*{arg\,max}_y$\\
\midrule
Operational CDT/source
 & $0$ & $9.3573397\cdot10^{-6}$ & $1.0019364869$ & $\{-378,378\}$\\
 & $1$ & $1.4950374\cdot10^{-6}$ & $1.0019292771$ & $\{-377,377\}$\\
\addlinespace
Literal PDF
 & $0$ & $0.0075371423$ & $1.1818109123$ & $\{-380,380\}$\\
 & $1$ & $0.0035191981$ & $1.1821799192$ & $\{-379,379\}$\\
\bottomrule
\end{tabular}
\end{table}

Thus Appendix~A's boundary-dominance claim fails under either interpretation.
The submitted source sets $M=1.002214572\ldots$.  This exceeds both required
\emph{operational} maxima in the first two rows, but it is below both maxima
under the literal PDF convention.  Every later statement that calls this
source value ``safe'' refers only to the operational CDT/source distribution.
The end-to-end experiments likewise use Equation~\eqref{eq:opgauss}; they do
not claim to instantiate and test a separate literal-PDF signer.

\subsection{Why the fixed order becomes a leakage channel}

With a valid envelope, the total assigned probability mass is at most one and
the cumulative intervals lie in $[0,1]$; the realized output law is then
independent of their scan order.  With $M_{\mathrm{bdry},c}$, central
candidates often contribute more than the entire available threshold range.
Under the central-minus order, the first candidate is $v=0$ for $c=0$ and
$v=-1$ for $c=1$.  Thus
\begin{equation}
                             v=-c+\delta,              \label{eq:delta}
\end{equation}
where the deviation $\delta$ is rare enough in the accepted transcript that
averaging exposes the secret.

Algorithm~3 says only to scan $V_c$ ``in a fixed order.''  No other part of
the submitted specification fixes that order, so the central-minus
permutation is a permitted deterministic choice fixed independently of the
key and target secret.  The violation of the stated envelope and the resulting
order dependence follow from Lemma~\ref{lem:envelope} and
Table~\ref{tab:envelope}.  The magnitude and correlation of the remaining
error after all signing gates are empirical properties of this completion.
Sections~\ref{sec:recovery} and \ref{sec:experiment} test those properties
using only signatures accepted and decoded by the pristine submitted
verifier.

\section{Public secret-tail recovery}
\label{sec:recovery}

For transcript $t$, let $C_t$ be the $256\times256$ negacyclic convolution
matrix of its public challenge polynomial $c_t$.  For bottom row $r$,
substitution of Equation~\eqref{eq:delta} into
Equation~\eqref{eq:response} gives the exact decomposition
\begin{equation}
 z_{t,r}=-C_t s_r+n_{t,r}+D_t s_r,                    \label{eq:model}
\end{equation}
where $n_{t,r}=y_{\mathrm{bot},t,r}$ and $D_t$ is convolution by
$\delta_t$.  Define the aggregate error
$\epsilon_{t,r}=n_{t,r}+D_t s_r$.  We do not assume that either summand, or
their sum after all signing gates, has exactly zero mean or is independent of
$C_t$.  Least squares is a heuristic dominant-term approximation whose
adequacy is assessed by public-key checks, held-out residuals, and the
end-to-end forgery experiment.

The ordinary least-squares estimate is
\begin{equation}
 \widehat{s}_r=-\left(\sum_t C_t^{\trans}C_t\right)^{-1}
                    \left(\sum_t C_t^{\trans}z_{t,r}\right).
                                                               \label{eq:ls}
\end{equation}
Negacyclic convolution diagonalizes at the odd $512$th roots of unity.  Let
$\widehat c_{t,j}$ and $\widehat z_{t,r,j}$ be the corresponding Fourier
coordinates.  Then Equation~\eqref{eq:ls} is evaluated componentwise as
\begin{equation}
 \widehat s_{r,j}=-
 \frac{\sum_t\overline{\widehat c_{t,j}}\widehat z_{t,r,j}}
      {\sum_t|\widehat c_{t,j}|^2}.                   \label{eq:fft}
\end{equation}
An inverse negacyclic FFT followed by coefficientwise rounding gives an
integer candidate $\widetilde s_{\mathrm{tail}}$.

The attack needs only values available from an ordinary signature:
\begin{enumerate}[leftmargin=1.7em]
  \item decode the signature response and derive $c_t$ from its transmitted
        challenge hash;
  \item append $(c_t,z_{\mathrm{bot},t})$ to the normal equations;
  \item at a public checkpoint, evaluate Equation~\eqref{eq:fft} and round;
  \item use Equation~\eqref{eq:keyrel} as an algebraic screen; and
  \item construct the fresh-message response, check its response and norm
        bounds and encodability, and accept the candidate only if the pristine
        verifier accepts the resulting signature.
\end{enumerate}
The public relation alone is not a certificate of forging capability: the
candidate's constructed response must also satisfy the verifier's bounds and
the submitted encoding.  The attack therefore stops only after both the
public relation and a fresh-message verification succeed.  A candidate has no
standing merely because its coefficients are small or its residual is
favorable.  Post hoc comparison with withheld truth determines whether the
accepted candidate is the original tail.  In the reported runs, the
specification-tape candidate first passed these tests at 40,000 signatures;
the independent-randomness candidate first passed at 30,000.

Equation~\eqref{eq:fft} can be accumulated online with $O(dn)$ complex state
and $O(Tdn\log n)$ arithmetic for $T$ signatures.  The reported reference
implementation instead uses a batch array for simplicity; its memory use is
reported in Section~\ref{sec:experiment}.

\section{Fresh-message forgery}
\label{sec:forgery}

The recovered tail gives a direct, deterministic forgery.  Let $M^*$ be a
message not submitted to the signing oracle and parse
$\pk=(\mathrm{seed}_A,b_{\mathrm{gen}})$.  At the specification level,
Algorithm~7 requires
\begin{equation}
 \mu^*=H_{\mathrm{gen}}(\mathrm{seed}_A,b_{\mathrm{gen}},M^*),\qquad
 c^*=\SampleInBall(0,\mu^*).                         \label{eq:specforgehash}
\end{equation}
For the submitted Rhyme-SHAKE-128 verifier, compute the exact byte-level
instantiation
\begin{align}
 \mu^*&=\SHAKE(\operatorname{enc}(\pk)\mathbin\Vert M^*)
                    [0\mathbin{:}\mathsf{CRHBYTES}],\nonumber\\
 \widetilde c^*&=\SHAKE(\operatorname{pack}_w(0)\mathbin\Vert\mu^*)
                    [0\mathbin{:}\mathsf{CTILDEBYTES}],\nonumber\\
 c^*&=\SampleChallenge(\widetilde c^*).              \label{eq:sourceforgehash}
\end{align}
This is the source implementation of
Equation~\eqref{eq:specforgehash}.  Set
\begin{equation}
             z^*=\bigl(c^*,s_{\mathrm{tail}}c^*\bigr)  \label{eq:forge}
\end{equation}
and call the submitted $\operatorname{pack\_sig}$ routine on
$(\widetilde c^*,z^*)$ to obtain the ordinary API signature.

\begin{theorem}[Zero-commitment forgery]
\label{thm:forge}
Let $s_{\mathrm{tail}}$ have coefficients in $[-2,2]$ and satisfy
Equation~\eqref{eq:keyrel}.  Then the signature in
Equation~\eqref{eq:forge} passes Rhyme's verification equation for $M^*$.
For the submitted Level-128 coefficient bounds, it also passes every response
norm bound.
\end{theorem}

\begin{proof}
Writing $s=(1,s_{\mathrm{tail}})$ and using commutativity of multiplication in
$\R$, the verifier reconstructs
\[
 Az^*-qjc^*=(As-qj)c^*=0\pmod {2q}.
\]
It therefore recomputes the zero commitment used to derive
$\widetilde c^*$.  More explicitly, the submitted verifier obtains
$c^*=\SampleChallenge(\widetilde c^*)$, reconstructs $w'=0$, and recomputes
\[
 \widetilde c^{\,\prime}=\SHAKE(\operatorname{pack}_w(0)\mathbin\Vert\mu^*)
                    [0\mathbin{:}\mathsf{CTILDEBYTES}]=\widetilde c^*.
\]
Thus its bytewise challenge-seed comparison succeeds.

The challenge has 30 coefficients equal to one and the rest zero.  Since
each coefficient of $s_{\mathrm{tail}}$ has magnitude at most two, every
coefficient of each product $s_rc^*$ has magnitude at most $60$.  Hence
\[
 \infnorm{z_1^*}\le1,\qquad
 \infnorm{z_{\mathrm{bot}}^*}\le60,\qquad
 \norm{z^*}_2^2\le30+4\cdot256\cdot60^2=3{,}686{,}430.
\]
These are below all three limits in Equation~\eqref{eq:bounds}.
\end{proof}

The construction does not recover $K$, regenerate the original serialized
secret key, search for a hash collision, or call the signing oracle after
recovery.  It forges each chosen target message directly.

\section{Machine-local end-to-end experiments}
\label{sec:experiment}

\subsection{Platform and source provenance}

All results in this section were generated on the same desktop: a four-core
Intel Core i5-6500 at 3.20~GHz with 16~GiB of RAM, running 64-bit Linux.  The
software environment used GCC~13.3.0, Python~3.13.13, and NumPy~2.5.3.  The
machine was shared with other audit jobs, so wall time varied with system load
and filesystem cache state.

The official candidate archive had SHA-256
\begin{center}
\small\texttt{7b509d21c6674bc23751b8743cb7c2a4a07ee295fafe04e2fa95c0613160bdfd},
\end{center}
and the specification had SHA-256
\begin{center}
\small\texttt{b1d6f86e7b2ce58a54e3078396c0df2283c59ae7773b3513580a19a8cb3c9ae3}.
\end{center}
All 49 files in the pristine Rhyme-SHAKE-128 tree used for external
verification were byte-identical to the corresponding files in that archive.

The tested signing completion made the following choices:
\begin{enumerate}[leftmargin=1.7em]
  \item retain the operational CDT/source distributions and use the two
        parity-specific values of $M_{\mathrm{bdry},c}$ obtained by applying
        Appendix~A's prescribed formula to those distributions;
  \item traverse the specified parity sets in the fixed central-minus order;
        and
  \item if the submitted first-response encoder cannot represent an otherwise
        accepted response, report encoding failure and let the signer's
        existing loop restart.
\end{enumerate}
The last rule addresses a separate source-level coverage gap without changing
an entropy frequency or the submitted decoding format.  The signer retained
the submitted entropy tables.  A preliminary extended-table variant was
discarded after its outputs failed pristine verification; none of its
measurements are used as attack-oracle evidence.  Consequently, every emitted oracle signature
could be passed to a verifier linked entirely against the untouched submitted
source.  The external verifier also decoded the response used to create each
public recovery record; the recovery experiment did not rely on a modified
decoder.

Algorithm~6 explicitly supplies the same tape $\mathit{rt2}$ to both the
sampling of $X'$ and \textsc{Sample Z}.  In the operational source mapping,
Run~1 therefore used the same stream identifier for the first component of
$X'$ and for rejection thresholds.  Run~2 changed only that rejection-stream
identifier to an unused domain.  This removes the direct random-stream reuse
between the thresholds and the first $X'$ component while retaining the
invalid boundary $M$, central-minus order, and all surrounding signing logic.
The two runs used distinct deterministic DRNG seeds and no shared state, so
their per-signature pseudorandom streams were separate across processes.
Deterministic seeds make the experiment
reproducible; the recovery process received neither seed and read only public
transcripts.

\subsection{Protocol}

For each attack run, a fresh key was generated and 50,000 distinct messages
were signed.  The pristine submitted verifier accepted and decoded
50,000/50,000 emitted attack-oracle signatures in each run, with zero
verification or decoding failures.  Recovery records were generated only from
those pristine-decoded signatures.  The first 40,000 records formed a common
comparison prefix and the last 10,000 were held out.  Before recovery, the
scoring secret was removed from the attack prefix.  Recovery and forgery then
ran with no secret file present; their outputs were hashed before the scoring
secret was restored for post hoc comparison.

The experimental collector generated all 50,000 signatures before running
the regression in order to reserve a held-out suffix.  The attack computations
used only the stated prefixes.  Because both the public relation and final
verification are public stopping tests, a prospective execution can run them
at the same fixed checkpoints and stop after 40,000 queries in Run~1 or
30,000 in Run~2; no held-out signature or target-secret diagnostic is needed.

The held-out residual is a diagnostic only.  It was not used for model
selection or stopping.  At public checkpoints the attack tested the rounded
tail against Equation~\eqref{eq:keyrel}, constructed the fresh-message
response, checked its bounds and encoding, and submitted it to the pristine
verifier.

\subsection{Recovery results}

Table~\ref{tab:recovery} reports post hoc coefficient scores.  Truth was used
only to populate this table; it was not an estimator input.  In Run~1, the
30,000-signature candidate failed Equation~\eqref{eq:keyrel}, while the
40,000-signature candidate passed the relation and forged.  In Run~2, the
30,000-signature candidate already passed and forged; its exactness persisted
at 40,000.

\begin{table}[t]
\centering
\caption{Exact rounded secret-tail coefficients out of 1,024.}
\label{tab:recovery}
\begin{tabular}{@{}rrrr@{}}
\toprule
Signatures & Run 1: reused tape & Run 2: independent tape & Paired safe-$M$\\
\midrule
10,000 & 995   & 989   & 380\\
20,000 & 1,021 & 1,022 & 400\\
30,000 & 1,023 & 1,024 & 399\\
40,000 & 1,024 & 1,024 & 399\\
\bottomrule
\end{tabular}
\end{table}

At 40,000 signatures, Run~1 had coefficient RMSE $0.120619$ and maximum
unrounded error $0.426050$.  Run~2 had RMSE $0.115143$ and maximum error
$0.382580$.  Their held-out residual RMS values were respectively
$124.218685$ and $124.661083$, compared with zero-model values $124.332177$
and $124.770305$.  No zero-mean premise is inferred from these measurements;
they show empirically that the $-C_ts_r$ term is recoverable by averaging for
the two tested keys.

Collection of all 50,000 signatures took 178.30 seconds for Run~1 and 160.44
seconds for Run~2, or 280.43 and 311.64 emitted, pristine-accepted signatures
per second.  The corresponding 40,000-signature batch recoveries took 9.87
and 6.97 seconds.  Peak recovery memory was 2.52~GiB in both runs.  The online
accumulation described after Equation~\eqref{eq:fft} removes this batch-memory
cost.

\subsection{Encoding-restart measurement}

The submitted first-response table has 96 high symbols and three raw low
bits.  With center $388$, it represents $z_1\in[-388,379]$ but has no escape
symbol for the verifier-allowed values $380,\ldots,388$.  We instrumented the
guard at that exact condition and separately counted every other
$\operatorname{pack\_sig}$ failure.  While producing 50,000 emitted signatures,
Run~1 restarted 5,696 times solely at this guard and Run~2 restarted 5,578
times.  These are respectively $10.23\%$ and $10.04\%$ of all candidates that
reached the packing stage, or 11.39 and 11.16 encoder-only restarts per 100
emitted signatures; neither run had another packing failure.  The gate is
therefore non-negligible and selects against the positive $z_1$ tail.  The
instrumented replays produced byte-identical public keys and oracle files to
the two reported attack collections.

To test whether that selection creates the recovery signal, we repeated Run~2
on the same key and messages with a diagnostic full-range first-response
table, so no otherwise accepted response was restarted by this gate.  This
table is incompatible with the pristine decoder, so the resulting signatures
are used only for this leakage ablation and are not counted as attack-oracle
signatures.  The ungated regression recovered 988, 1,021, 1,024, and 1,024
exact coefficients at 10,000, 20,000, 30,000, and 40,000 signatures.  At
30,000 its coefficient RMSE was $0.134043$, compared with $0.131398$ for the
pristine-encodable Run~2 transcript; both tails were already exact.  The
ungated recovered tail also satisfied the public relation and produced a
pristine-verifier-accepted fresh-message forgery.  Thus the restart gate
slightly changes sample selection but does not account for the recovery
channel in this paired experiment.

\subsection{Forgeries and controls}

The recovered tails matched all 1,024 separately retained scoring
coefficients byte for byte.  The pristine verifier accepted a fresh-message
forgery for each run:
\begin{center}
\begin{tabular}{@{}lrrr@{}}
\toprule
Run & Signature bytes & Squared norm & Pristine verification\\
\midrule
1 & 1,372 & 28,464 & accept\\
2 & 1,372 & 28,636 & accept\\
\bottomrule
\end{tabular}
\end{center}
The signing-query messages contained a fixed prefix followed by a counter;
the 28-byte forgery message was absent from both query sets.  Three controls
were run for each key: the forgery under another message, a signature formed
after changing one recovered coefficient, and the forgery under an
independently generated public key.  All six checks rejected.

For Run~1, the public key, 50,000-record pristine-decoded transcript,
recovered tail, and pristine-codec forgery have SHA-256 values
\begin{center}
\small
\texttt{18cb64be10afa7d869870c62bd32c9b314af9078711dd9cb9882dbe0c9ce2246},\\
\texttt{29abbf892d1706907fe1caf6fbeebba61c96856e1c402cba2baaf9a4ef70d0bf},\\
\texttt{01e3e36399fc89aa7c806e75cad5146abdd67c324384c51cfe793652719746b7},\\
\texttt{38fb4b7ecfa6939504de183287fa596b6cc45c84c69b1f049f637f64bffa154f}.
\end{center}
For Run~2, the corresponding values are
\begin{center}
\small
\texttt{a1d5d5b475e275baabfc894acbb42064575efeb843b3170e7fbf5003b40cc6b2},\\
\texttt{d97c47f98e2cb8a0c0cb4fae225ce5f720e60df97ca5bc52d00c4b163ebc7dba},\\
\texttt{de72d6d1e48747e7aa5d06dbc25a82bd8a1a5ef1384a3423f4202044065375a1},\\
\texttt{8f663a36111e357ca5f7a7d106f81a051362d3dc6801f6f2a4ea0af72ff72113}.
\end{center}

\subsection{Paired safe-envelope control}

The safe-$M$ control was rerun with exactly the same key and ordered messages
as Run~2; the public-key files are byte-identical.  It retained the
central-minus order, independent rejection stream, submitted entropy tables,
encoding-failure restart, and all surrounding signing logic.  A source diff
between the two instrumented signers contains one behavior-changing edit:
the sampler uses the source's $M=1.002214572\ldots$ instead of the two
operational boundary values.  The pristine verifier accepted and decoded
50,000/50,000 emitted control signatures.

The paired control incurred 2,603 encoder-only restarts while producing those
50,000 signatures, or $4.95\%$ of the 52,603 candidates reaching the packing
stage; it had no other packing failure.  Applying the same regression to the
first 40,000 signatures produced an all-zero rounded candidate.  Only 399 of
its 1,024 coefficients happened to match the secret; its coefficient RMSE was
$0.971998$ and its maximum unrounded error was $2.318174$.  The held-out
residual RMS and zero-model RMS were both $129.564098$ to the displayed
precision.  The public relation failed and the attempted forgery was rejected.
This same-key, one-variable ablation is consistent with correcting $M$
removing the particular first-moment channel exploited here.  It remains one
empirical control and does not prove the corrected implementation secure
against every possible attack.

\section{Scope, limitations, and repair}
\label{sec:scope}

The result gives an end-to-end EUF-CMA adversary against the tested
Rhyme-SHAKE-128 completion: using a prefix of at most 40,000 chosen-message
signatures, it outputs an accepted signature for a message outside the full
50,000-message experimental query set.  Its measured time, memory, and query
costs are practical on the tested four-core desktop.  The experimental
collector's held-out suffix is excluded from the adversary's input, as
described in Section~\ref{sec:experiment}.

The following qualifications delimit that statement.
\begin{itemize}[leftmargin=1.5em]
  \item The submitted specification does not define a unique scan order.
        It states the correct envelope condition but subsequently prescribes
        a boundary closed form that violates it.  The attack targets the
        deterministic central-first completion described above.
  \item The unmodified submitted SHAKE source is not the vulnerable signer.
        It chooses ascending arrays and uses an envelope that is safe for its
        operational CDT distributions.  The same numerical value is not an
        envelope under the literal PDF Gaussian convention.
  \item The submitted first-response encoder does not cover every value
        allowed by the verification bounds.  The tested completion treats an
        unencodable response as a signing restart, without modifying the
        submitted entropy tables.  Every reported oracle signature was
        accepted and decoded by pristine submitted code.  This gate is
        non-negligible, but the same-key ungated diagnostic recovered at the
        same 30,000-signature checkpoint, so the gate was not needed for
        recovery in that comparison.
  \item End-to-end recovery was tested at the 128-bit parameter set on two
        independent keys.  Two successes do not estimate a key-averaged
        success probability or establish sample complexity at higher levels.
  \item The practical signing experiment uses the operational Gaussian
        distributions encoded by the submitted CDT and fixed-point tables.
        The literal PDF convention gives the separate values in
        Table~\ref{tab:envelope}.  Its boundary formula is also invalid, but
        no end-to-end experiment is claimed for a literal-PDF signer.
  \item The attack recovers the exact secret tail, not the original seed or
        every byte of the serialized signing key.
\end{itemize}

The specification should replace Equation~\eqref{eq:mbdry} by the exact
maximum in Equation~\eqref{eq:envelope} or by a proved upper bound.  It should
test the resulting cumulative mass, specify a canonical enumeration order,
domain-separate every random tape, define one Gaussian parameter convention,
and require lossless encoding for every response accepted by verification.
The security proof and parameter tables should then be recomputed for that
single corrected sampler.

\section{Conclusion}

The submitted specification states the correct cumulative-envelope
requirement, but then derives and prescribes a boundary closed form that
violates that requirement for the submitted Rhyme-128 parameters.  Together
with the unspecified fixed enumeration order, the invalid constant makes the
sampler order-dependent.  In the tested central-first operational completion,
the public response became a recoverable noisy linear observation of the
secret tail.  On this machine, two independent Level-128
tails were recovered from at most 40,000 pristine-accepted signatures and
yielded accepted fresh-message forgeries through the pristine verifier.
The same-key control using the operationally valid source envelope was
consistent with removal of this regression channel, and an ungated diagnostic
showed that the encoder restart was not needed for recovery in the paired
comparison.  These results
establish a practical break of the tested completion and a
specification/design-level defect in an internally inconsistent
submitted specification.  They do not claim that every implementation of the
underlying abstract Rhyme construction is broken; in particular, the
submitted source's operationally safe envelope lies outside this attack.

\begin{thebibliography}{9}

\bibitem{rhyme-spec}
Zhongxiang Zheng, Anyu Wang, Chunhuan Zhao, Guangwu Xu, Sibo Feng,
Zhichen Yan, Yuhang Sun, Zhengtao Jiang, Congming Wei, Shuang Sun,
Chen Jia, and Shuaigang Li.
\newblock \emph{RHYME: Algorithm Specifications and Supporting Documentation}.
\newblock Submission sign-22 to the Next-generation Commercial Cryptographic
Algorithms Program, 2026.

\bibitem{ngcc-rhyme}
ngcc.dev.
\newblock \emph{Rhyme (sign-22): public security findings}.
\newblock \url{https://ngcc.dev/reports/sign-22.html}, accessed
1 October 2026.

\end{thebibliography}

\end{document}
